\chapter*{Preface: A Word About Knobs}
\addcontentsline{toc}{chapter}{Preface: A Word About Knobs}
\markboth{Preface}{}

Every theory of the universe comes with a control panel.

Some panels are tidy: a few dials, each labelled, each set once and then left alone. Some panels look like the cockpit of an airliner. And at least one, if its admirers are right, has a single dial with so many settings that you could not write them all down if every atom in the observable universe were a sheet of paper.

This book is about those panels. It is about the dials on them, which physicists call \emph{free parameters} and which I will mostly call \emph{knobs}, because that is what they are. A knob is a number the theory cannot tell you. You have to look at the sky, read off the answer, and turn the knob until the theory agrees. Every theory has at least one. The interesting questions are how many, which ones, and what happens when you turn them.

The book grew out of an experiment that I ran with an AI research team in the autumn of 2026 and wrote up under a deliberately silly title: \emph{ACDM flavored fried chicken version of Galaxy Dynamics}. The idea was simple and a little mischievous. I had a galaxy law, one formula with one coefficient, that describes the rotation of disc galaxies remarkably well. The standard model of cosmology, \LCDM, explains the same galaxies with halos of dark matter shaped by supernova explosions, black holes and the slow churn of galaxy formation. Those processes come with knobs, and lots of them. So I asked: if I let \LCDM\ turn every knob it has, how close can it get to my one formula? How many knobs does it take? And is there anything the knobs \emph{can't} do?

The answers were surprising, and they cut both ways. That is why this book exists.

\section*{What you will find here}

\textbf{Part I, The Knob Shop}, explains what a knob is, mathematically. It covers degrees of freedom, information criteria, Bayesian evidence and the Occam factor that quietly charges you for every dial you add. Then it takes a respectful tour of the six knobs that run \LCDM's fit to the cosmic microwave background, which are genuinely among the great achievements of science, and an increasingly less respectful tour of the dozens of knobs that appear once you try to make a galaxy.

\textbf{Part II, The Fried Chicken Kitchen}, is the experiment itself. Twelve supernova knobs. Fourteen. Twenty-six. A recipe with 1026 parameters, one with about 1066, and a halo defined directly from my formula that matches it to four decimal places, because the formula went in. You'll see exactly how far knobs get you, which is very far, and the one ingredient that no knob could supply, which is the law's tie to dark energy.

\textbf{Part III, A Knob With $10^{500}$ Clicks}, takes a friendly but pointed look at the string theory landscape. Its famous count of vacua makes the fried chicken kitchen look like a lemonade stand.

\textbf{Part IV, What About One Knob?}, asks the obvious next question. What if a theory had one knob, $\kappa=\tfrac12$, and nothing else at galaxy scale? I lay out what that one knob gets right in galaxies, and then, just as carefully, where it fails. It fails in clusters, in the cosmic microwave background, in a few stubborn dwarf galaxies, in the vertical structure of the Milky Way and elsewhere. A book about counting knobs is worthless if it doesn't count its own.

\textbf{Part V, The Knob That Can't Be Turned}, ends with the part I care about most: a prediction with no knob left in it. Once $\kappa$ is set by galaxies today, the law says how its acceleration scale must change across cosmic time, following the dark energy density measured by surveys that have nothing to do with galaxy rotation. Nobody can tune that, not me and not anyone else. It will either come out right or it won't.

\section*{How to read it}

The book is meant to be fun and technical at the same time. There are jokes, a running cooking metaphor and some chapter openers in a voice I'd never use in a journal. There are also equations, real ones, with every symbol defined, and tables of numbers from frozen, committed calculations. If you like the equations, you'll find derivations in the blue \emph{math boxes}. If you don't, skip them; the prose will carry you. The red \emph{knob boxes} count the knobs in whatever we're discussing. The dark \emph{honesty boxes} say where a claim is weaker than it sounds. Please read those even if you skip everything else.

A unit you'll meet constantly is the \emph{dex}: a difference of 1 dex is a factor of 10, and 0.1 dex is a factor of about 1.26. When I say a recipe sits ``0.0415 dex from the law,'' I mean the typical gap between them, root-mean-square over thousands of points, is a factor of $10^{0.0415}\approx1.10$, about ten per cent.

\section*{A promise}

I'll say this once here and keep saying it: \textbf{nothing in this book shows that my law is right and \LCDM\ is wrong.} The fried chicken experiment shows that \LCDM's knobs can imitate the law closely. That is evidence of \emph{freedom}, not of \emph{failure}. My law fails in places where \LCDM\ succeeds beautifully, above all in the cosmic microwave background. Counting knobs is a way of asking how much a theory has stuck its neck out, not a way of declaring a winner.

What a single unturnable knob can do is let the universe declare the winner. The last chapter is about that.

\bigskip
\noindent\emph{Briar Creek, November 2026}
